\section{Complete finite odd-$C_2$ endpoints}

\paragraph{Status.}
D-OC2-F is a complete author derivation with an actually executed exact
certificate, awaiting independent whole-proof review. Its limited coverage
is explicit: the full universal cover at $q=5$, and center-positive
representations only at $q=7,9$. No generic odd-$C_2$ table is used.

\begin{lemma}[D-OC2-F1: exact full-normalizer endpoint bounds]
Let $V$ be a nontrivial irreducible complex representation of one of
$\operatorname{Sp}_4(5)$, $\operatorname{PSp}_4(7)$ or
$\operatorname{PSp}_4(9)$. Every finite-order nonscalar operator
normalizing its image has projective support at least $\dim(V)/4$.
The minimum degrees are respectively $12,25,41$.
\end{lemma}

\begin{proof}
The retained complete ordinary table input consists of
\[
\begin{gathered}
\texttt{2.S4(5)},\quad\texttt{2.S4(5).2},\qquad
\texttt{S4(7)},\quad\texttt{S4(7).2},\\
\texttt{S4(9)},\quad\texttt{S4(9).2\_1},\quad
\texttt{S4(9).2\_2},\quad\texttt{S4(9).2\_3}.
\end{gathered}
These are official complete CTblLib tables \cite{OC2:CTbl}. The two $q=5$ tables have
ATLAS provenance. The $q=7,9$ tables are computations of Thomas Breuer
using Dixon's algorithm and character-theoretic methods, dated 2001
and 2004 in the retained source strings; they are not mislabeled as
ATLAS computations. The full character values, class sizes and orders,
prime power maps, normal fusions, and provenance strings are retained.
Their input SHA-256 is
\texttt{6824ca705e8619733ced3c7aceab6e145b7b45735fa56a9990ed7607eb6f7a73}.

The GAP exporter verifies full row orthogonality. A separate exact
cyclotomic-polynomial program selects every irreducible core row and,
in an extension, every row restricting irreducibly to that core. It
reconstructs every eigenvalue multiplicity by the Fourier formula
\[
       m_j(x)=\frac1d\sum_{k=0}^{d-1}
                  \chi(x^k)\exp(-2\pi i j k/d),\qquad d=|x|.
\]
Scalar entries are identified by $|\chi(x)|^2=\chi(1)^2$. Exact
arithmetic on all 485 selected rows and 37502 nonscalar entries gives
no violation of the quarter bound. Normal class unions, class orders,
group orders and fusion compatibility with all retained prime power
maps are checked as well. The minimum support ratios over the
respective core and extension tables are
\[
                     4/13,\qquad9/25,\qquad16/41.
\]
The minimum nontrivial core degrees are $12,25,41$.
The complete verifier and its previously reviewed cyclotomic helper are
retained; no floating-point approximation enters these assertions.

We explain why these tables cover the claimed normalizers. The outer
groups at $q=5,7$ are cyclic of order two, and at $q=9$ they are
$C_2\times C_2$, by the symplectic row of
\cite[Section~2.2.2, Table~2.1, author p.~14]{HES:LiXia2022}.
The latter three cyclic subgroups are exactly the field, diagonal and
diagonal-field extensions identified in the official source strings.
Every individual operator belongs, projectively, to one of these
cyclic extension cases; a representation of the whole outer group is
not needed. At $q=5$ the multiplier-two universal cover and its outer
extension are both retained, so center-negative as well as
center-positive modules are covered.

More explicitly, an intertwiner for a compatible automorphism is unique
up to scalar by Schur's lemma. Choose an element of the listed extension
inducing the same automorphism, after multiplying by a core element if
necessary. Its cyclic power relation differs from the intertwiner's
power by a scalar; rescaling the intertwiner solves that relation over
$\mathbf C$. Thus it supplies an irreducible extension character whose
restriction is the given core module. Changing the cyclic lift or its
central power relation changes only scalar factors in the spectra.
This is why one listed covering extension suffices also for the
isoclinic alternative at $q=5$. All selected rows were scanned, so the
operator has the asserted support. A core representation factoring
through its central quotient remains covered by the same input.
\end{proof}

\begin{theorem}[D-OC2-F2]
An inducing image in the signature-$(3,3)$ sixfold-source problem with
an irreducible normal core formed from the covered groups in D-OC2-F1
cannot have an actual period image of dimension nine. Arbitrary
compatible scalar/outer extensions and permutations of tensor factors
are allowed. The claim for the $q=7,9$ factors is restricted to
center-positive modules.
\end{theorem}

\begin{proof}
The core representation is a tensor product of its factor modules.
For a factor-preserving operator, normalized eigenvalue multiplicity
distributions convolve, and their maximum is at most that of any
nonscalar factor. A nontrivial cycle of length $\ell$ on degree-$a$
factors is a weighted permutation of $a^\ell$ basis words, with $a$
constant words and all other orbits of length at least two. Its largest
eigenspace is at most $a+(a^\ell-a)/2\le3a^\ell/4$.
Thus all nonscalar operators in the full normalizer have support at
least $e/4$, where $e$ is the total degree. A nonidentity scalar has
actual moving rank $e$.

The actual intermediate-cover variation has
$6=2(h-1)e+\sum_j\operatorname{rank}(g_j-I)$.
Since $e\ge12$, genus $h\ge2$ is impossible; genus one has at most
six effective branch points and hence actual period-image dimension
at most six. A nine-dimensional image would require genus zero and
at least twelve branch points. Its moving-rank sum $2e+6$ would then
be at least $12e/4=3e$, forcing $e\le6$, a contradiction.
No self-duality, determinant or equality-of-character-fields assumption
is used in this argument.
\end{proof}

\paragraph{Exact remaining coverage gate.}
The executed inventory found no stored complete universal-cover tables
for $\operatorname{Sp}_4(7)$ or $\operatorname{Sp}_4(9)$, and no stored
complete odd-$C_2$ tables at $q=11,13,17,19,23,25,27$.
The preliminary full generic odd-$C_2$ table is known to contain errors,
and the available generic unipotent table is not a complete ordinary
table \cite{OC2:GCT}. Neither is used or extrapolated here. The already reviewed Weil
module and large-field exclusions retain their separate scopes.
